\chapter{Foi, initiation et Église}\label{ch:foi_eglise}

\section{Néophytes et catéchumène}

Le vocabulaire de l’initiation est rare mais net. Sept défunts sont dits \emph{neofytus} ou \emph{neofyta}, mot toujours lu sans restitution : Romanus, mort à 9 ans en 371 (\nnum{145}) ; Libera, à 8 ans en 374 (\nnum{172}) ; Quiminatus, à environ 5 ans en 379 (\nnum{210}) ; Victor, dont l’âge est perdu, en 382 (\nnum{233}) ; Leontius, à environ 28 ans en 385 (\nnum{276}) ; Aristo, à 8 mois en 389 (\nnum{308}) ; Aufinianus, à 42 ans en 396 (\nnum{391}). Un défunt, dont le nom est transmis sous la forme \emph{Bonifae}, est dit \emph{catechuminus} et meurt à environ 60 ans en 397 (\nnum{405}).

Les âges des néophytes vont de huit mois à quarante-deux ans : la mention \emph{neophytus} signale un baptême récent, non un âge. Quatre des sept néophytes sont des enfants. Les dates de décès ou de déposition conservées ne se concentrent pas dans le temps pascal (5 mai 374, 24 novembre 379, 14 août 382, 22 juin 385, 4 juin 389) : la mention indique un statut au moment de la mort, non une circonstance liturgique. Elle peut toutefois signaler un baptême reçu à l’approche de la mort, ce que suggèrent le très jeune âge d’Aristo et l’association, au \nnum{308}, de \emph{neofytus} et de \emph{puer innox}.

Le \nnum{405} est le seul texte du corpus qui nomme un catéchumène. Le défunt est mort à environ soixante ans sans avoir été baptisé, ou du moins sans que l’épitaphe le dise. Le texte n’en est pas moins chrétien par sa forme ; il atteste que le statut de catéchumène pouvait être gravé, et donc revendiqué, sur une tombe.

\section{Les fidèles et la foi}

Huit textes emploient \emph{fides} ou \emph{fidelis}. Dans quatre d’entre eux, l’adjectif \emph{fidelis} qualifie le défunt comme membre de la communauté des baptisés : Secunda, \emph{virgo fidelis} (\nnum{70}, 362) ; Celerinus, enfant d’un an et quatre mois, \emph{fidelis quiescis in pace} (\nnum{231}, 382) ; un garçon de dix ans (\nnum{278}, 383) ; Aufinianus, \emph{ratione fidelis} et néophyte (\nnum{391}, 396). Le cas de Celerinus suggère que l’enfant était baptisé ; celui d’Aufinianus montre que les deux qualifications pouvaient être cumulées.

Dans d’autres textes, \emph{fidelis} et \emph{fides} relèvent de l’éloge moral ou conjugal : l’épouse d’Alexander n’a jamais connu de discorde avec lui, \emph{qua fidelis} (\nnum{115}) ; Secunda a vécu \emph{pura fide} (\nnum{70}) ; la défunte du \nnum{382} a montré une grande fidélité (\emph{tante fides}). Le \nnum{133} qualifie Felix de \emph{sanctae fidei vocitus}, « appelé à la sainte foi », expression qui peut évoquer une conversion ou un baptême, sans que le texte précise. L’épitaphe de Maximilla (\nnum{306}) emploie trois fois le mot \emph{fides} au sens de la foi chrétienne : \emph{propositumque fidem} (vers 1), \emph{si animo recolat fidem datam} (vers 2), \emph{mirae fidei} (dans la partie en prose).

Un seul texte emploie \emph{christianus} : la défunte du \nnum{201} (377) est dite \emph{[- - -] Christianae uncta [- - -]}, « ointe chrétienne » ou « ointe de la [foi ?] chrétienne », selon la restitution du contexte, qui est perdu. La lecture du mot est assurée dans les deux bases ; sa construction ne l’est pas.

\section{L’Église et le clergé}

Les mentions de l’Église sont toutes liées à des fonctions cléricales, et toutes sont en partie restituées : \emph{lect[or] ec(c)les[iae sanctae c]at(h)olic(a)e} (\nnum{67}, 362) ; \emph{[p]r(es)b(yter) [eccl(esiae) cathol(icae)]} (\nnum{428}, 398) ; \emph{[ecclesiae ?] catholica[e]} (\nnum{18}, 346), où l’EDB conserve le point d’interrogation de l’éditeur. Seul l’adjectif \emph{catholica} est lu au \nnum{18}, et en partie au \nnum{67} ; le mot \emph{ecclesia} n’est lu en entier dans aucun texte. L’adjectif \emph{catholica} qualifie l’Église comme institution ; on ne peut en inférer une intention polémique.

Les autres mentions de clercs — lecteurs de titres, prêtres, diacres, évêques — sont étudiées au chapitre~\ref{ch:rangs}. Les titres de la Ville apparaissent à travers leurs lecteurs : \emph{titulus Fasciolae} (\nnums{202} et 421), \emph{Pudentiana} (\nnum{269}).

\section{Virginité et ascèse}

Le mot \emph{virgo} est lu dans six textes (\nnums{70}, 99, 183, 306, 425 et 458). Il peut désigner une jeune fille non mariée ou une vierge consacrée, et le contexte ne permet pas toujours de choisir. La consécration est probable lorsque le mot est associé à \emph{ancilla dei} (Maximilla, \nnum{306} ; Praetiosa, morte à douze ans, \nnum{458}) ou à \emph{fidelis} (Secunda, \nnum{70}). Elle ne l’est pas pour Petronia, \emph{virgo} morte à dix-neuf ans (\nnum{183}). L’éloge de Quintilianus (\nnum{474}), qui « aime la chasteté » et « rejette le monde », appartient au même registre ascétique, de même que la \emph{castitas} de la défunte du \nnum{382} et de Severa, épouse du prêtre Gaudentius (\nnum{310}).

\section{Bilan}

Le vocabulaire de la foi et de l’initiation (\emph{fidelis}, \emph{neophytus}, \emph{catechumenus}, \emph{ancilla dei}, \emph{virgo}) est plus fréquent que celui de la doctrine. Il qualifie la personne et sa place dans la communauté. Le mot \emph{baptisma} n’apparaît pas ; le baptême est désigné par ses effets, le statut de néophyte ou de fidèle.
